% ============================================================================
% GÉNÉRÉ par scripts/exemples-guide.ts — NE PAS ÉDITER À LA MAIN.
% Régénérer : npm run exemples (chaque chiffre est vérifié contre le moteur).
% ============================================================================
\pexemples La prestation maximale dépend de l'âge des enfants ; la réduction, de
l'$\AFNI$ (tableau~\ref{tab:bases-exemples}) et du \emph{nombre} d'enfants
(taux par palier).

\medskip\noindent\textbf{M6 --- famille monoparentale, 35~ans, revenu de travail de \D{35000}, un enfant de 3~ans en garderie subventionnée (\D{2000} payés dans l'année).}
Un enfant de moins de 6~ans (maximum \num{7997}). L'$\AFNI$ de
\num{32685} — abaissé par la déduction pour frais de garde (poste~6) — est
\emph{sous} le premier seuil de \num{37487} : aucune réduction.
\begin{align*}
ACE &= \num{7997} - 0 = \num{7997}
\end{align*}
Prestation \emph{maximale} : \D{7997}, non imposable.

\medskip\noindent\textbf{M8 --- couple, 38 et 36~ans, revenus de travail de \D{60000} et \D{40000}, deux enfants : 4~ans (garde non subventionnée, \D{13000}) et 8~ans (garde non subventionnée, \D{3000}).}
Deux enfants (4 et 8~ans) : maximum de \num{7997} $+$
\num{6748} $=$ \num{14745} ; taux de réduction « 2~enfants »
(\pc{13.5} puis \pc{5.7}).
\begin{align*}
\text{bande 1} &= \min\bigl(\AFNI - \num{37487},\ \num{43735}\bigr) = \num{43735} \\
\text{bande 2} &= \pos{\num{86070} - \num{81222}} = \num{4848} \\
\text{réduction} &= \pc{13.5} \times \num{43735} + \pc{5.7} \times \num{4848} = \num{6180.56} \\
ACE &= \num{14745} - \num{6180.56} = \num{8564.44}
\end{align*}
L'$\AFNI$ (\num{86070}) bénéficie ici de la déduction fédérale pour frais
de garde de \num{13000} (poste~6) : sans elle, la réduction serait plus
forte. ACE : \D{8564.44}.

\medskip\noindent\textbf{M9 --- couple, 45 et 45~ans, revenus de travail de \D{120000} et \D{60000}, un enfant de 5~ans (garde non subventionnée, \D{15000}).}
Un enfant de moins de 6~ans, $\AFNI$ de \num{170361} (taux
« 1~enfant » : \pc{7} puis \pc{3.2}).
\begin{align*}
\text{réduction} &= \pc{7} \times \num{43735} + \pc{3.2} \times \num{89139} = \num{5913.90} \\
ACE &= \pos{\num{7997} - \num{5913.90}} = \num{2083.10}
\end{align*}
Même à ce niveau de revenu, il reste \D{2083.10} : le second palier ne
réduit qu'à \pc{3.2}.
